% Generated draft: requires mathematical and source review.
\documentclass[11pt]{article}
\usepackage[margin=25mm]{geometry}
\usepackage{fontspec}
\setmainfont{Latin Modern Roman}
\usepackage{amsmath,amssymb,mathtools,unicode-math}
\setmathfont{Latin Modern Math}
\usepackage{enumitem,needspace,xurl,hyperref}
\setlength{\parindent}{0pt}
\setlength{\parskip}{5pt}
\setlength{\emergencystretch}{3em}
\newcommand{\N}{\mathbb{N}}
\newcommand{\Z}{\mathbb{Z}}
\newcommand{\Q}{\mathbb{Q}}
\newcommand{\R}{\mathbb{R}}
\newcommand{\C}{\mathbb{C}}
\newcommand{\F}{\mathbb{F}}
\newcommand{\A}{\mathbb{A}}
\newcommand{\PP}{\mathbb{P}}
\newcommand{\E}{\mathbb{E}}
\newcommand{\eps}{\varepsilon}
\newcommand{\dd}{\,\mathrm{d}}
\begin{document}
\section*{UnsolvedMath ID: 30004536}
\textbf{Source problem code:} catalogue entry
\subsection{Definitions and mathematical statement}
The supplied record poses a family of questions about integrals of self-similar Markov processes, but does not specify enough hypotheses to determine a unique mathematical statement. The following notation describes a conditional formulation; the choices introduced here are not asserted to be those intended by the source.

For a real-valued Markov family \((Y,\mathbb P_y)_{y\in E}\), where \(E\subseteq\R\) is invariant under positive dilations, self-similarity with index \(H>0\) means that, for every \(c>0\) and \(y\in E\),
\[
 (Y_{ct})_{t\geq0}\text{ under }\mathbb P_y
 \ \stackrel{\mathrm{law}}{=}\ 
 (c^H Y_t)_{t\geq0}\text{ under }\mathbb P_{c^{-H}y}.
\]
This definition presupposes that the processes are defined for the indicated times; any killing, lifetime, or boundary convention must also be specified. If \(Y\) has locally time-integrable paths, define its integral, with initial integral value \(z\in\R\), by
\[
 I_t=z+\int_0^t Y_s\,\dd s.
\]
The initial data for the integrated process generally include both \(I_0=z\) and \(Y_0=y\); write \(\mathbb P_{z,y}\) for such a joint initial condition when defined.

For example, if the intended exit problem concerns a specified finite interval \((a,b)\), with \(a<z<b\), one could define
\[
 \tau_{(a,b)}=\inf\{t\geq0:I_t\notin(a,b)\},
 \qquad \inf\varnothing=\infty.
\]
If \(I\) is continuous and defined for all relevant times, and if \(T_u=\inf\{t\geq0:I_t=u\}\), then \(\tau_{(a,b)}=T_a\wedge T_b\). The source does not define \(T_a\), \(T_b\), or the subscript \(x\), so this interpretation remains conditional.

\textbf{Requested problems.} Determine persistence and exit-time asymptotics for the integral of a general self-similar Markov process. In particular, the supplied expression to be studied is
\[
 \mathbb P_x(T_a\wedge T_b>t),\qquad t\longrightarrow\infty.
\]
Under the interval interpretation above, this would ask for the large-time behavior of \(\mathbb P_{z,y}(\tau_{(a,b)}>t)\), for the specified process, boundary conditions, and admissible initial data. The source does not state whether an exponent, an asymptotic equivalent, a leading constant, or another description is required, and does not separately define the persistence event.

The record also asks for the analogous exit-time problem when the driving process \(Y\) is a stable L\'evy process, and asks to compute the area under a general stable excursion. If \(e\) denotes a specified excursion with lifetime \(\zeta\), its signed area, when the integral exists, would be
\[
 A(e)=\int_0^{\zeta}e(s)\,\dd s.
\]
Neither the underlying stable process, the excursion law and normalization, nor the meaning of ``compute'' is supplied. In particular, the source does not determine whether this is a signed or nonnegative area, or whether the desired output is its distribution, transform, moments, or asymptotics. These qualifications must be resolved before these questions become fully specified mathematical problems.
\subsection{Short English statement}
The problem asks for long-time survival probabilities for integrated self-similar Markov processes, including stable L\'evy processes, and for the area law of a general stable excursion. The source leaves the process class, initial data, boundary conventions, and excursion normalization unspecified. The displayed definitions therefore give only a conditional framework requiring human mathematical review.
\subsection{Sources}
\begin{itemize}
\item Supplied local UnsolvedMath registry record, permanent ID \(30004536\), source problem code \texttt{catalogue entry}, entitled ``Persistence and Exit Times of Integrated Self-Similar Processes,'' in the Oberwolfach Reports Open Problems collection. No individual proposer is supplied.
\item Supplied report citation: \emph{Stochastic Processes under Constraints (2021)}, Oberwolfach Reports, DOI \texttt{10.4171/owr/2020/32}, pp.\ 1601--1656. \url{https://doi.org/10.4171/owr/2020/32}. The registry gives proposed year \(2020\); the supplied citation uses \(2021\). This dating has not been independently checked.
\item Supplied literature reference labeled ``Profeta--Simon'': \url{https://doi.org/10.1007/s00440-014-0577-5}.
\item Supplied literature reference labeled ``Stable excursion area'': \url{https://arxiv.org/abs/2009.01504}.
\item Literature assessment checked at \texttt{2026-08-22}, according to the supplied dated curation, not independently verified here. The supplied status is ``partially solved.'' The assessment states that important one-sided and spectrally positive cases of exit and persistence questions for integrated stable processes have been solved, including excursion-area results, while the general stable excursion and all listed variants remain open. The source expressly calls this dated literature triage derived from the supplied OWR curation and requests independent verification of time-sensitive status and cited claims.
\item The registry labels the statement \texttt{corrected\_verified}, with the qualification that only formula delimiters, notation, grammar, and typography were repaired and that mathematical content and scope were preserved. This supplied label does not establish mathematical correctness or independent source verification.
\end{itemize}
\end{document}
